% ============================================================
%  Black-Litterman with Ledoit-Wolf Shrinkage
%  Fixed-Omega leakage and the critical correlation threshold
%  Standalone paper draft - Lucas Posern
%
%  All quoted numbers are reproduced by code/paper_numbers.py from
%  the frozen data in ../data/ (excess returns, sample cov ddof=1,
%  sklearn LW intensity applied to the ddof-1 matrix, implied delta,
%  weights via (delta Sigma)^-1 mu).
%  Revision 2026-07: corrected rho* (factor-2 bug), exact per-pair
%  Delta c_eff table, updated backtest to thesis view set, paper
%  frame (abstract, introduction, conclusion), decimal points.
%  Revision 2026-07 v2 (referee fixes): single-view scoping of the
%  Idzorek result + Remark 1 (K>1), signed rho* statement + negative
%  branch, ddof convention pinned, weights = (delta Sigma)^-1 mu
%  throughout, frozen-Omega backtest column, Memmel test, corrected
%  mu_s inequality in 6.6, attribution fixes (Walters, He-Litterman),
%  new refs (Markowitz, Sharpe, Chopra-Ziemba, Memmel, LW 2008).
%  Extension pass (this file, fixed_omega.tex; formerly main_v2.tex): general view vectors
%  (criterion v vs mu_s*||p||^2, absolute/basket examples), range and
%  deflation floor (Remark 2), PD-boundary notes, worst-case c*;
%  numbers reproduced by section [7] of code/paper_numbers.py.
%  Revision 2026-07 v3 (post-review fixes): linearity claim corrected
%  (prior exact, projection first-order), Delta-pi interpretation
%  rewritten (covariance-with-portfolio form), view-channel scoping of
%  the Omega question, omega- vs c-primitive caveat in 6.8, long-only
%  twin figures reported (incl. frozen), accuracy claims 0.01/0.2pp,
%  deflation floor 1.05pp, alpha-hat = 0.0406 pinned, inline proof of
%  the Idzorek closed-form equivalence, language pass, 11 new refs.
%  Revision 2026-07 v4: out-of-sample backtest section commented out
%  (\iffalse block; the closed forms are identities, the weight- and
%  exposure-level consequences are estimation-window facts now reported
%  in the precision-channel section); Memmel/LW2008 refs parked;
%  strict style pass (em-dash density, active voice, contrast formulas);
%  "training window" -> "estimation window" throughout.
% ============================================================

\documentclass[a4paper,11pt]{article}

\usepackage[a4paper, left=2.5cm, right=2.5cm, top=2.5cm, bottom=2.5cm]{geometry}
\usepackage{amsmath, amssymb, amsthm, amsfonts}
\usepackage{mathtools}
\usepackage[english]{babel}
\usepackage[T1]{fontenc}
\usepackage[utf8]{inputenc}
\usepackage{booktabs}
\usepackage{array}
\usepackage{float}
\usepackage{graphicx}
\usepackage{caption}
\usepackage{xcolor}
\usepackage[round]{natbib}
\usepackage{listings}
\usepackage[colorlinks=true, linkcolor=blue!70!black, citecolor=blue!70!black, urlcolor=blue!70!black]{hyperref}

\lstset{
  language=Python,
  basicstyle=\ttfamily\footnotesize,
  keywordstyle=\color{blue!70!black}\bfseries,
  commentstyle=\color{green!40!black}\itshape,
  stringstyle=\color{red!70!black},
  breaklines=true,
  columns=fullflexible,
  showstringspaces=false,
  frame=single,
  framerule=0.4pt,
  rulecolor=\color{gray!60},
  xleftmargin=4pt,
  xrightmargin=4pt,
}

% ------------------------------------------------------------
% Theorem environments
% ------------------------------------------------------------
\theoremstyle{plain}
\newtheorem{proposition}{Proposition}
\newtheorem{theorem}[proposition]{Theorem}
\newtheorem{corollary}[proposition]{Corollary}
\theoremstyle{definition}
\newtheorem{example}[proposition]{Example}
\newtheorem{remark}[proposition]{Remark}

% ------------------------------------------------------------
% Notation
% ------------------------------------------------------------
\newcommand{\R}{\mathbb{R}}
\newcommand{\E}{\mathbb{E}}
\newcommand{\Sig}{\Sigma}
\newcommand{\Om}{\Omega}
\newcommand{\muBL}{\mu_{\mathrm{BL}}}
\newcommand{\Std}{\mathrm{Std}}
\newcommand{\LW}{\mathrm{LW}}
\newcommand{\eff}{\mathrm{eff}}

\title{Black--Litterman with Ledoit--Wolf Shrinkage:\\
Fixed-$\Omega$ Leakage and the Critical Correlation Threshold%
\thanks{This paper grew out of the author's bachelor thesis at TU Dresden.
I thank Rodrigo S.~Targino (FGV EMAp, Rio de Janeiro) for detailed feedback
on the proof strategy. All remaining errors are mine.}}
\author{Lucas Posern\\
\small Faculty of Mathematics, TU Dresden\\
\small \texttt{posernlucas@gmail.com}}
\date{\today}

% ------------------------------------------------------------
\begin{document}

\maketitle

\begin{abstract}
\noindent
The Black--Litterman model is often combined with shrinkage estimators
of the covariance matrix, most prominently the Ledoit--Wolf estimator.
We show that the effect of the estimator switch on the posterior depends
on how the view-uncertainty matrix $\Omega$ is handled. Under
Idzorek's confidence calibration, $\Omega$ co-moves with $\Sigma$ and, for
a single view, the switch is exactly inert in the view channel; with
several views the deviation is driven by the cross-view covariances
and, in our three-view application, stays an order of magnitude below
the effects studied here. If instead
$\Omega$ is calibrated once on the sample covariance and then frozen, a
workflow that can easily arise when the covariance model is upgraded after
the views have been signed off, the effective view confidence
$c_{\eff}$ shifts by a closed-form amount that is independent of the prior
scale $\tau$ and grows with the shrinkage intensity. The shift is positive
precisely when the view variance lies below twice the average asset
variance, a condition met by relative views on highly correlated pairs
\emph{and} on low-variance (bond) pairs; for a general view vector the
threshold is $\mu_s\lVert p\rVert^2$, and on our data an absolute view
on the lowest-variance asset shifts by $+12.9$ percentage points, more
than any pair. We derive an exact expression for
the induced changes in posterior means and weights under a single relative
view, provide a closed-form critical correlation threshold, and quantify
the effect on an ETF implementation of Idzorek's eight-asset universe: at
a shrinkage intensity of $0.041$, five of 28 pairs shift their effective
confidence by more than five percentage points, the largest by $+10.7$.
Recomputing $\Omega$ alongside $\Sigma$ removes the leakage exactly.
\end{abstract}

\medskip
\noindent\textbf{Keywords:} Black--Litterman; Ledoit--Wolf shrinkage;
view uncertainty; covariance estimation; portfolio construction.\\
\textbf{JEL:} G11; C13; C58.

% ============================================================
\section{Introduction}
\label{sec:intro}
% ============================================================

The Black--Litterman model \citep{BlackLitterman1992} is among the most
widely implemented models of quantitative portfolio construction. It
anchors expected returns at the CAPM market equilibrium
\citep{Sharpe1964} and blends in subjective views with an explicit
uncertainty matrix $\Omega$, thereby taming the input sensitivity of
mean--variance optimisation \citep{Markowitz1952} documented by
\citet{Michaud1989}, \citet{BestGrauer1991} and \citet{ChopraZiemba1993};
Bayesian shrinkage of the mean \citep{Jorion1986,FrostSavarino1986} and
decision-theoretic treatments of parameter uncertainty
\citep{KanZhou2007} attack the same problem from the estimation side.
Its second input, the covariance matrix $\Sigma$, receives comparatively
little attention in the Black--Litterman literature (empirical
exceptions include \citealp{BesslerOpferWolff2017}), although the sample
estimator is unreliable whenever the number of assets is not small
relative to the sample size. A standard remedy is the shrinkage
estimator of \citet{LedoitWolf2004}, which pulls the sample covariance
toward a scaled identity and improves the conditioning of $\Sigma^{-1}$
substantially: on our data the condition number falls by a factor of
three, and by far more as $n$ approaches $T$ \citep{LedoitWolfHoney2004}.

This paper studies what happens inside the Black--Litterman model when the
covariance estimator is switched from the sample estimator to Ledoit--Wolf.
Within the view channel, the answer turns on a bookkeeping question
that, to our knowledge, has not been isolated in the literature: the
treatment of $\Omega$. Two workflows arise naturally. In the first, $\Omega$ is
derived from $\Sigma$ through Idzorek's confidence calibration
\citep{Idzorek2007} and therefore co-moves with the estimator. In the
second, $\omega_k$ is computed once on the sample covariance, stored as a
set of scalars, and left untouched when the covariance model is later
upgraded. We call the resulting shift of the effective view confidence
away from the stated confidence, at frozen $\Omega$, \emph{leakage}.

Our contribution is threefold. First, we show that under Idzorek's
calibration the estimator switch is \emph{exactly} inert in the view
channel for a single view: the effective view absorption equals the
specified confidence $c$ for any positive-definite estimator, and the
remaining effect on the posterior runs through the equilibrium prior,
which is exactly linear in the shrinkage intensity $\alpha$, and the
view projection, which is linear in $\alpha$ to first order. With several
simultaneous views the deviation is driven by the cross-view covariances;
Remark~\ref{rem:multiview} quantifies it. Second, for the fixed-$\Omega$
workflow we derive a closed-form expression for the shift
$\Delta c_{\eff}$ of the effective view confidence. The expression is
exact, independent of $\tau$ (the same $\tau$ enters calibration and
posterior, and cancels), linear in the shrinkage intensity $\alpha$ for
small $\alpha$, and its sign is determined by a simple criterion: the view variance versus twice the
average asset variance. In the symmetric case of equal asset volatilities
this criterion translates into a closed-form critical correlation
threshold $\rho^*(\alpha, \varepsilon)$. An exact identity for the
Black--Litterman posterior under a single relative view then propagates
$\Delta c_{\eff}$ analytically into posterior means and portfolio
weights. Third, we quantify the effect on an ETF implementation of the
eight asset classes in \citet{Idzorek2007}: with the analytic shrinkage
intensity $\alpha = 0.041$, five of the 28 asset pairs shift their
effective confidence by more than five percentage points; among them is
a bond pair with moderate correlation, which shows that high correlation
is not necessary for a material shift. Separating the view channel from
the precision channel through $\Sigma^{-1}$ shows that a frozen
$\Omega$ restores roughly $40\,\%$ of the gross-exposure reduction that
the estimator switch delivers.

Covariance
estimator, confidence and view uncertainty form one unit; upgrading one
piece without re-deriving the others moves the effective model away from
the analyst's stated inputs. Related observations (that constraints or
estimators interact with estimation error in ways invisible at the input
stage) appear in \citet{JagannathanMa2003} for long-only constraints and
in \citet{DeMiguelGarlappiUppal2009} for the difficulty of beating naive
diversification out of sample. Within the Black--Litterman literature,
\citet{HeLitterman1999}, \citet{SatchellScowcroft2000},
\citet{Herold2003} and \citet{Walters2014} document the model's
sensitivity to the $(\tau, \Omega)$ specification (see
\citealp{Meucci2010} for a survey, \citealp{KolmRitter2017} for the
Bayesian structure, \citealp{MichaudEschMichaud2013} for a critical
view, and \citealp{BevanWinkelmann1998} for a production workflow); our
result gives one mechanism of this sensitivity a closed form. On the
shrinkage side, \citet{LedoitWolf2003}, the nonlinear extensions of
\citet{LedoitWolf2017} and \citet{LedoitWolf2020}, and the
portfolio-oriented intensity calibration of
\citet{DeMiguelMartinUtreraNogales2013} would enter our formulas only
through a different expression for the view-variance shift; the leakage
mechanism itself is target-independent.

Section~\ref{sec:blm-brief} fixes notation,
model and data. Sections~\ref{sec:sample-problem} and~\ref{sec:lw} review
the conditioning problem and the Ledoit--Wolf estimator.
Section~\ref{sec:idzorek-cancel} settles the Idzorek case.
Section~\ref{sec:fixed-omega} contains the main results for fixed
$\Omega$. Section~\ref{sec:precision} isolates the precision channel
through $\Sigma^{-1}$ and quantifies the weight-level cost of the
frozen-$\Omega$ workflow; Section~\ref{sec:conclusion} concludes.

% ============================================================
\section{The Black--Litterman Model and the Data}
\label{sec:blm-brief}
% ============================================================

For $n$ risky assets with covariance $\Sigma\in\R^{n\times n}$,
$\Sigma\succ 0$, the Black--Litterman posterior mean is
\begin{equation}
\label{eq:bl-posterior}
  \muBL \;=\; \bigl[(\tau\Sig)^{-1} + P^\top\Om^{-1}P\bigr]^{-1}
              \bigl[(\tau\Sig)^{-1}\pi \;+\; P^\top\Om^{-1}Q\bigr],
\end{equation}
with implied equilibrium return $\pi = \delta\Sigma w_{\mathrm{eq}}$,
view-pick matrix $P\in\R^{K\times n}$, view returns $Q\in\R^K$,
view-uncertainty matrix $\Om\succ 0$, prior scale $\tau>0$, and risk
aversion $\delta>0$. The unconstrained portfolio is
$w_{\mathrm{BL}} = \tfrac{1}{\delta}\Sigma^{-1}\muBL$.%
\footnote{Some implementations map the posterior to weights with
$\Sigma_{\mathrm{post}} = \Sigma + M$, where $M$ is the posterior
covariance of the mean \citep{HeLitterman1999}. All results and all
reported numbers in this paper use $\Sigma$ itself; for relative views
this keeps the weights exactly fully invested.}

\paragraph{Data.}
We use an ETF implementation of the eight asset classes in
\citet{Idzorek2007}: AGG (US bonds), BWX (international bonds), IWF/IWD
(US large-cap growth/value), IWO/IWN (US small-cap growth/value), EFA
(international developed equity) and EEM (emerging markets equity). We
take monthly total returns from Yahoo Finance (adjusted closes) and
convert them to excess returns with the 13-week T-bill rate; the series
are complete over the whole analysis period. All estimates use the
window January 2008 to December 2018 ($T=132$, $n=8$). The market weights
$w_{\mathrm{eq}}$ are the market-capitalisation weights published by
Idzorek for these asset classes; they predate the sample (2004 vintage),
which makes them a dated but look-ahead-free prior anchor. All data,
the estimation pipeline and every number quoted in this paper are
reproduced by a single script; see the reproducibility statement at the
end.

\paragraph{Calibrated quantities.}
On the estimation window the implied risk aversion of the S\&P~500 excess
return series is $\delta = 2.44$, the ratio of mean to variance of
monthly excess returns (the annualisation factor cancels); the
S\&P~500 serves as the proxy for $\delta$ while the equilibrium weights
come from Idzorek's asset classes, a mismatch we accept as the
customary approximation. We use $\tau = 0.05$, the upper end of the range in
\citet{HeLitterman1999} and \citet{Idzorek2007} (the central result of
Section~\ref{sec:fixed-omega} is independent of this choice), and view
confidences $c_k = 0.50$ throughout, an agnostic midpoint; the
$c$-dependence of the leakage is explicit in the factor $c(1-c)$ of
\eqref{eq:dceff-exact}. The average asset variance, which
acts as the Ledoit--Wolf shrinkage target scale, is
$\mu_s = \tfrac{1}{n}\mathrm{tr}(\hat\Sigma_s) = 0.0282$ annualised
($2.35\cdot 10^{-3}$ per month), and the analytic shrinkage intensity is
$\hat\alpha = 0.0406$, reported as $\alpha = 0.041$ throughout. Three
relative views with these confidences run through
the multi-view sections: IWF $>$ IWD by $+3\,\%$ p.a., EEM $>$ EFA by
$+4\,\%$ p.a., and IWF $>$ AGG by $+6\,\%$ p.a. The single-view analysis
of Section~\ref{sec:fixed-omega} uses the first of these.

% ============================================================
\section{The Sample Covariance Problem}
\label{sec:sample-problem}
% ============================================================

The sample covariance estimator is
\[
  \hat\Sigma_s \;=\; \frac{1}{T-1}\sum_{t=1}^T (r_t - \bar r)(r_t - \bar r)^\top.
\]
For $n$ assets and $T$ observations its expected estimation loss grows
with $n/T$ \citep{LedoitWolf2004}. Once $n/T$ is not
negligible, the eigenvalues of $\hat\Sigma_s$ spread
too widely; large eigenvalues are overestimated, small ones are
underestimated \citep[p.~371]{LedoitWolf2004}. The smallest eigenvalue
moves toward zero, the condition number
$\kappa(\Sigma) = \lambda_{\max}/\lambda_{\min}$ grows, and the inverse
$\hat\Sigma_s^{-1}$ used in
$w_{\mathrm{BL}} = \tfrac{1}{\delta}\Sigma^{-1}\muBL$ amplifies
estimation error. On the estimation window the condition number is
$\kappa(\hat\Sigma_s) = 341.8$.

% ============================================================
\section{Ledoit--Wolf Shrinkage}
\label{sec:lw}
% ============================================================

\citet{LedoitWolf2004} regularise $\hat\Sigma_s$ by shrinking it toward
the scaled identity $\mu_s I_n$, with
$\mu_s = \tfrac{1}{n}\mathrm{tr}(\hat\Sigma_s)$,
\[
  \hat\Sigma_{\LW} \;=\; (1-\alpha)\,\hat\Sigma_s \;+\; \alpha\,\mu_s\,I_n,
\]
where the oracle coefficient $\alpha\in[0,1]$ minimises
$\E[\|(1-\alpha)\hat\Sigma_s + \alpha\mu I_n - \Sigma\|_F^2]$, with
$\mu$ the population counterpart of $\mu_s$; the analytic estimator
replaces the unobservable moments by consistent plug-ins, yielding
$\hat\alpha$. Because $\alpha\mu_s I_n$ adds $\alpha\mu_s>0$ to
every eigenvalue of $(1-\alpha)\hat\Sigma_s$, $\hat\Sigma_{\LW}\succ 0$
is automatic. Throughout, we estimate $\hat\alpha$ with the analytic
implementation of \texttt{scikit-learn} and apply it to the
$T-1$-normalised sample matrix, so that the identity above holds exactly
in the empirical pipeline. (The covariance output of
\texttt{scikit-learn} itself normalises by $T$ and would rescale every
closed form below by $(T-1)/T$.) On the estimation window the analytic
estimator yields $\hat\alpha = 0.0406$ and brings the condition number down
to $\kappa(\hat\Sigma_{\LW}) = 108.0$, a factor $3.2$ reduction.

% ============================================================
\section{Idzorek $\Omega$: a Single View Is Inert}
\label{sec:idzorek-cancel}
% ============================================================

When $\Omega$ is derived from $\Sigma$ via
Idzorek's calibration, the estimator cancels exactly in the view channel
of a single view.

Idzorek's confidence calibration leads, in the closed form derived by
\citet{Walters2014}, to
\[
  \omega_{kk} \;=\; \frac{1-c_k}{c_k}\,\tau\,p_k^\top\Sigma p_k
\]
(Idzorek's original procedure matches weight tilts iteratively; under
the weight mapping of Section~\ref{sec:blm-brief} the view tilt equals
$c_{\eff}$ times the full-confidence tilt by \eqref{eq:bl-single-w}, so
matching a fraction $c$ of it is solved exactly by this $\omega_{kk}$
in the single-view case). For a single view $p$
the posterior then simplifies along the view direction to
\[
  p^\top\muBL \;=\; (1 - c)\,p^\top\pi \;+\; c\,Q,
\]
an identity we derive in Section~\ref{sec:dmu-dw}. The view absorption
equals $c$ exactly, independent of $\Sigma$. Switching
$\hat\Sigma_s \to \hat\Sigma_{\LW}$ scales the view variance and the
matching $\omega$ proportionally, the ratio stays at $c$. Along the view
direction the remaining LW impact runs through the prior spread
$p^\top\pi$ alone; off the view direction the projection
$\Sigma p / v_k$ moves as well. The prior shift is exactly linear in
$\alpha$; the projection is a smooth rational function of $\alpha$,
linear to first order (and exactly linear only when $p$ is an
eigenvector of $\hat\Sigma_s$).

Concretely, the LW estimator shifts the prior linearly,
\[
  \Delta\pi_i \;=\; \alpha\,\delta\,\Bigl[(\mu_s - \sigma_i^2)\,w_{\mathrm{eq},i}
        \;-\; \sum_{j\neq i}\hat\Sigma_{s,ij}\,w_{\mathrm{eq},j}\Bigr].
\]
Writing $\sigma_i^2 := \hat\Sigma_{s,ii}$, the diagonal term raises the
prior of below-average-variance assets, but the cross-covariance term
dominates in equity-heavy universes: the shift collapses to
$\Delta\pi_i = \alpha\,\delta\,[\mu_s\,w_{\mathrm{eq},i} -
\mathrm{Cov}(r_i, r_p)]$ with $r_p = w_{\mathrm{eq}}^\top r$, so assets
that co-move strongly with the equilibrium portfolio lose prior return.
On our universe all six equity ETFs have $\Delta\pi_i < 0$ (including
IWF and IWD, whose variances lie below $\mu_s$), the two bond ETFs
gain, and $|\Delta\pi_i| \leq 0.22\,\%$ p.a.\ throughout.

\begin{remark}[several views]\label{rem:multiview}
For $K > 1$ the absorption matrix is $A = S\,(S + \Om)^{-1}$ with
$S = \tau P\Sigma P^\top$. Idzorek's $\Omega$ is tied to
$\mathrm{diag}(S)$ alone (up to the factor $(1-c_k)/c_k$), so
off-diagonal elements of $S$ break both the exact
absorption $c$ and its estimator-invariance. On the three-view set of
Section~\ref{sec:blm-brief} the residuals
$P\muBL - [(1-c)P\pi + cQ]$ stay below $0.6\,\%$ p.a., and the per-view
absorption remains within one percentage point of $c = 50\,\%$ and moves
by less than $0.3$ percentage points between the two estimators, an
order of magnitude below the fixed-$\Omega$ shifts of
Section~\ref{sec:fixed-omega}.
\end{remark}

This settles the Idzorek case for a single view; the fixed-$\Omega$ case
does not cancel.

% ============================================================
\section{Fixed $\Omega$: the Leakage Mechanism}
\label{sec:fixed-omega}
% ============================================================

A second workflow computes
$\omega_k$ once on the sample estimator and stores the scalars. Later,
the analyst replaces the covariance with an LW estimator, perhaps
because the badly conditioned $\hat\Sigma_s^{-1}$ produced implausible
leverage, but leaves $\omega_k$ fixed. We are not aware of survey evidence on how
widespread this workflow is, though interfaces that accept a
user-supplied $\Omega$ while $\Sigma$ is estimated elsewhere make it
easy to adopt; the point of this section is that following it changes
the model's effective inputs without any visible change at the input
stage. The view absorption shifts implicitly: the view variance $v_k$
(defined below) moves with the estimator while $\omega_k$ stays fixed.

We derive the closed form for the shift in $c_{\eff}$, translate it into
a general sign-and-size criterion and, for the symmetric case, into a
critical correlation threshold, and verify the chain on the IWF/IWD
pair. Throughout this section the analysis rests on a single view, taken as
the relative pair $p = e_i - e_j$ unless stated otherwise; the
symmetric per-component statements additionally assume
$\sigma_i \approx \sigma_j$ for the view pair, an assumption we make
explicit where it is used and drop where it is not needed.

\subsection{Effective View Absorption $c_{\eff}$}
\label{sec:ceff}

The actual view absorption after evaluating the posterior is
\[
  c_{\eff,k} \;=\; \frac{\tau\,v_k}{\tau\,v_k + \omega_k},
  \qquad v_k := p_k^\top\Sigma p_k.
\]
$c_{\eff,k}$ measures how far the posterior travels along the view
direction from prior toward view. At $c_{\eff,k} = 0$ the posterior
stays at the prior, at $c_{\eff,k} = 1$ it lands exactly on the view.
The ratio depends on the view variance $v_k$, the view uncertainty
$\omega_k$, and the prior scale $\tau$, which controls how tightly the
prior is held: at fixed $\omega_k$, small $\tau$ pulls $c_{\eff,k}$
down, large $\tau$ lets the posterior move toward the view more easily.

In the Idzorek world with $\omega_k = \tfrac{1-c_k}{c_k}\tau v_k$, both
$\tau$ and $v_k$ cancel and $c_{\eff,k} = c_k$ exactly. In the
fixed-$\Omega$ world this cancellation breaks.

\subsection{Shift of $v_k$ Under Shrinkage}
\label{sec:vk-shift}

For a relative view $p_k = e_i - e_j$ the view variance is
$v_k = \sigma_i^2 + \sigma_j^2 - 2\hat\Sigma_{s,ij}
     = \sigma_i^2 + \sigma_j^2 - 2\rho_{ij}\sigma_i\sigma_j$,
with no assumption on the volatilities. Under LW,
\[
  v_k^{\LW} \;=\; (1 - \alpha)\,v_k^\Std \;+\; 2\alpha\,\mu_s,
\]
since $p_k^\top I_n p_k = 2$. Hence $v_k$ grows under shrinkage exactly
when $v_k^\Std < 2\mu_s$, i.e.\ when the view variance lies below twice
the average asset variance. For positively correlated pairs of similar
volatility this is the typical case; it also holds for pairs whose
variances are simply small relative to the universe average, such as
bond pairs. At fixed $\omega_k$, the absorption $c_{\eff,k}$ grows with
$v_k$.

The pair structure of the view enters only through $p_k^\top I_n p_k$.
For an arbitrary view vector $p_k$,
\[
  v_k^{\LW} \;=\; (1-\alpha)\,v_k^\Std \;+\; \alpha\,\mu_s\,\lVert p_k\rVert^2,
\]
and every formula below holds verbatim with $2\mu_s$ replaced by
$\mu_s\lVert p_k\rVert^2$: the general sign criterion is
$v_k^\Std \lessgtr \mu_s\lVert p_k\rVert^2$. Two special cases follow
immediately. An \emph{absolute} view on asset $i$ ($p_k = e_i$,
$\lVert p_k\rVert^2 = 1$) inflates exactly when $\sigma_i^2 < \mu_s$:
on our universe a frozen-$\omega$ absolute view on AGG shifts by
$+12.9$ percentage points, the largest leakage in the universe, ahead
of every pair, while an absolute view on EEM
($\sigma_{\mathrm{EEM}}^2 > \mu_s$) shifts by $-0.5$ percentage
points, the deflation regime realised in the data. A basket view
(equal-weighted US growth over US value, $\lVert p_k\rVert^2 = 1$)
shifts by $+5.8$ percentage points.

\subsection{Closed Form for $\Delta c_{\eff,k}$}
\label{sec:dceff}

Combining $c_{\eff,k}(\hat\Sigma_{\LW})$ and $c_{\eff,k}(\hat\Sigma_s)$
over a common denominator and simplifying,
\begin{align*}
  \Delta c_{\eff,k}
  \;&=\; \frac{\tau\,v_k^\LW}{\tau\,v_k^\LW + \omega_k}
        \;-\; \frac{\tau\,v_k^\Std}{\tau\,v_k^\Std + \omega_k}\\
  \;&=\; \frac{\tau\,v_k^\LW\,(\tau\,v_k^\Std + \omega_k) \;-\; \tau\,v_k^\Std\,(\tau\,v_k^\LW + \omega_k)}
              {(\tau\,v_k^\LW + \omega_k)(\tau\,v_k^\Std + \omega_k)}\\
  \;&=\; \frac{\tau\,\omega_k\,(v_k^\LW - v_k^\Std)}
              {(\tau\,v_k^\LW + \omega_k)(\tau\,v_k^\Std + \omega_k)}.
\end{align*}
Substituting $v_k^\LW - v_k^\Std = \alpha(2\mu_s - v_k^\Std)$ yields
\begin{equation}
\label{eq:dceff}
  \Delta c_{\eff,k}
  \;=\;
  \frac{\tau\,\omega_k\,\alpha\,(2\mu_s - v_k^\Std)}
       {(\tau\,v_k^\LW + \omega_k)(\tau\,v_k^\Std + \omega_k)}.
\end{equation}
For small $\alpha$ the shift is linear in $\alpha$, and its sign matches
$\mathrm{sgn}(2\mu_s - v_k^\Std)$ for any $\alpha$: the general criterion
announced above, valid for arbitrary volatility ratios.

In the workflow under study, $\omega_k$ was calibrated on the sample
estimator, $\omega_k = \tfrac{1-c}{c}\,\tau\,v_k^\Std$. Substituting
this into \eqref{eq:dceff} cancels $\tau$ completely and gives the exact
form
\begin{equation}
\label{eq:dceff-exact}
  \Delta c_{\eff,k}
  \;=\;
  \frac{c\,(1-c)\,\bigl(v_k^\LW - v_k^\Std\bigr)}
       {c\,v_k^\LW + (1-c)\,v_k^\Std}.
\end{equation}
Three consequences are worth recording. First, the leakage does not
depend on the prior scale $\tau$ at all; no choice of $\tau$ mitigates
it. Second, \eqref{eq:dceff-exact} is exact for any pair of positive
view variances, so it can be evaluated asset pair by asset pair without
any symmetry assumption. Third, at fixed variances the shift is
maximised over the confidence at
$c^* = 1/\bigl(1+\sqrt{v_k^{\LW}/v_k^\Std}\,\bigr)$; on the IWF/IWD
pair $c^* = 0.45$ with a shift of $+10.9$ percentage points, so the
agnostic $c = 0.50$ used throughout operates close to the worst case.
As $v_k^\Std \to 0$ at fixed $\mu_s$ the shift approaches its supremum
$1-c$.

\begin{remark}[range and boundary]\label{rem:range}
As a function of $v_k^\Std \in (0,\infty)$ at fixed target
$\mu_s\lVert p_k\rVert^2$, the shift \eqref{eq:dceff-exact} is strictly
decreasing (the numerator of its derivative is
$-c(1-c)\,\alpha\mu_s\lVert p_k\rVert^2$), with range
\[
  \Delta c_{\eff,k} \;\in\;
  \Bigl(-\,\frac{c\,(1-c)\,\alpha}{1-c\alpha},\; 1-c\Bigr).
\]
The leakage is therefore asymmetric: the inflation end is bounded only
by $1-c$ (fifty percentage points at $c = 0.50$), while the deflation
end has an $O(\alpha)$ floor: at $\alpha = 0.041$ and $c = 0.50$ no
view in any universe can lose more than $1.05$ percentage points of
effective confidence, and even at $\rho_{ij} = -1$ in the symmetric
setting of Section~\ref{sec:rho-star} the shift is $-0.4$
($\alpha = 0.041$) to $-3.4$ percentage points ($\alpha = 0.30$). Both
endpoints are boundary limits: $v_k^\Std \to 0$ corresponds to
$\rho_{ij} \to 1$ for the equal-volatility pair view, and
$\rho_{ij} = \pm 1$ itself lies on the boundary of the
positive-definite cone, where the posterior machinery degenerates but
\eqref{eq:dceff-exact}, which needs only $v_k^\Std > 0$, extends
continuously. Note also that shrinkage enforces
$v_k^{\LW} \geq \alpha\mu_s\lVert p_k\rVert^2 > 0$: Ledoit--Wolf
regularises exactly the near-degenerate view directions in which the
frozen-$\omega$ confidence explodes.
\end{remark}

\subsection{Critical Correlation $\rho^*$ in the Symmetric Case}
\label{sec:rho-star}

To turn \eqref{eq:dceff-exact} into a rule of thumb, consider the
symmetric case $\sigma_i = \sigma_j = \sigma$, in which
$v_k^\Std = 2\sigma^2(1 - \rho_{ij})$. Setting
$\Delta c_{\eff,k} = \varepsilon$ in \eqref{eq:dceff-exact} and solving
for $\rho_{ij}$ with $v_k^\LW = (1-\alpha)v_k^\Std + 2\alpha\mu_s$
yields the closed-form threshold
\begin{equation}
\label{eq:rho-star}
  \rho^*(\alpha,\varepsilon)
  \;=\;
  1 \;-\; \frac{c\,\alpha\,\mu_s\,[(1-c) - \varepsilon]}
              {\sigma^2\,[\varepsilon\,(1 - c\,\alpha) \;+\; c\,(1-c)\,\alpha]},
  \qquad \varepsilon \in (0,\,1-c).
\end{equation}
A view triggers $\Delta c_{\eff,k} > \varepsilon$ exactly when
$\rho_{ij} > \rho^*$. The threshold falls as $\alpha$ rises and as the
tolerated shift $\varepsilon$ falls. (On the negative branch,
$\Delta c_{\eff,k} < -\varepsilon$ requires $\rho_{ij}$ below a second
threshold $\rho^{**}$, which exists once
$\varepsilon < c(1-c)\alpha/(1-c\alpha)$; at $\alpha = 0.30$ and
$\varepsilon = 1\,\%$ it sits at $\rho^{**} = -0.32$, at
$\alpha = 0.041$ it lies far outside $[-1,1]$. The tables below concern
the positive branch; Remark~\ref{rem:range} bounds the entire negative
branch.) The symmetric case serves as a rule of
thumb: for $\sigma_i \neq \sigma_j$
the exact criterion remains \eqref{eq:dceff-exact} evaluated at the
heterogeneous $v_k^\Std = \sigma_i^2 + \sigma_j^2 -
2\rho_{ij}\sigma_i\sigma_j$, and we use that exact form for the
pair-level results below.

\subsection{Values and Economic Statement}
\label{sec:values}

Table~\ref{tab:rho-star} shows $\rho^*$ for four shrinkage intensities
and four tolerance levels $\varepsilon$. The parameters $\mu_s = 4.0\cdot
10^{-3}$, $\sigma^2 = 3.5\cdot 10^{-3}$, $c = 0.50$ are illustrative,
chosen to match the order of magnitude of the monthly quantities in our data
($\mathrm{tr}(\hat\Sigma_s)/n = 2.35\cdot 10^{-3}$ per month); the table
depends on the parameters only through the ratio $\mu_s/\sigma^2$
and $c$.

\begin{table}[H]
\centering
\caption{Critical correlation $\rho^*(\alpha,\varepsilon)$ from
\eqref{eq:rho-star}. The pair correlation $\rho_{ij}$ must exceed this
value for $\Delta c_{\eff} > \varepsilon$. Illustrative parameters
$\mu_s = 4.0\cdot 10^{-3}$, $\sigma^2 = 3.5\cdot 10^{-3}$, $c = 0.50$.}
\label{tab:rho-star}
\begin{tabular}{lcccc}
\toprule
$\varepsilon$ & $\alpha = 0.04$ & $\alpha = 0.10$ & $\alpha = 0.166$ & $\alpha = 0.30$ \\
\midrule
$1\,\%$  & $0.434$ & $0.188$ & $0.083$ & $-0.006$ \\
$2\,\%$  & $0.629$ & $0.377$ & $0.239$ & $0.106$ \\
$5\,\%$  & $0.826$ & $0.645$ & $0.511$ & $0.343$ \\
$10\,\%$ & $0.915$ & $0.810$ & $0.715$ & $0.571$ \\
\bottomrule
\end{tabular}
\end{table}

At moderate shrinkage the threshold is demanding only for tight
tolerances: with $\alpha = 0.04$, a five-percentage-point shift already
occurs above $\rho^* = 0.826$, a correlation level typical of style
pairs within one equity market ($\rho = 0.920$ for IWF/IWD,
Table~\ref{tab:pairs}). At strong shrinkage ($\alpha = 0.30$),
every non-negatively correlated pair crosses the one-percent band. The
analytic intensity tends to grow with $n/T$ \citep{LedoitWolf2004}, so
larger universes on the same window typically sit further to the right
in the table.

The symmetric threshold understates the reach of the effect, because the
general criterion is $v_k^\Std < 2\mu_s$, not high correlation per se.
Table~\ref{tab:pairs} evaluates the exact expression
\eqref{eq:dceff-exact} for every pair of the universe, one view at a
time, with the full $8\times 8$ covariance shrunk once at
$\alpha = 0.041$ (the practitioner workflow) and $c = 0.50$; the table
reports the pairs above the five-point band. Five of 28 pairs shift by
more than five percentage points. The top of the list is dominated by
the highly correlated, nearly homogeneous style pairs, led by IWF/IWD at
$+10.7$ percentage points: the model operates at an effective
confidence of $0.61$ where the analyst specified $0.50$. But the bond
pair AGG/BWX ranks third at $+8.2$ percentage points despite a moderate
correlation of $0.57$: its view variance is small relative to $2\mu_s$
because bond variances are far below the universe average. A threshold
stated in correlations alone would miss it. In this universe every one
of the 28 pairs satisfies $v_k^\Std < 2\mu_s$, so the leakage is a
one-sided confidence inflation; the deflation regime
$v_k^\Std > 2\mu_s$ requires pairs whose spread variance exceeds twice
the universe average, such as weakly correlated high-volatility assets.

\begin{table}[H]
\centering
\caption{Exact $\Delta c_{\eff}$ from \eqref{eq:dceff-exact} for all
pairs with a shift above five percentage points; each pair is evaluated
as a single relative view in isolation. Full-universe shrinkage at
$\alpha = 0.041$, $\omega$ frozen on the sample estimator,
$c = 0.50$; $k = \sigma_i/\sigma_j$. Estimation window 2008--2018.}
\label{tab:pairs}
\begin{tabular}{lccc}
\toprule
Pair & $\rho_{ij}$ & $k$ & $\Delta c_{\eff}$ (pp) \\
\midrule
IWF/IWD & $0.920$ & $0.99$ & $+10.7$ \\
IWO/IWN & $0.939$ & $1.03$ & $+8.8$ \\
AGG/BWX & $0.574$ & $0.45$ & $+8.2$ \\
IWD/IWN & $0.917$ & $0.80$ & $+6.6$ \\
IWF/IWO & $0.909$ & $0.77$ & $+5.5$ \\
\bottomrule
\end{tabular}
\end{table}

In economic terms, whenever a relative view connects two
assets whose spread variance is small (because they are highly
correlated, or because both carry little variance to begin with), a
frozen $\Omega$ lets the shrinkage-induced widening of the spread
variance flow into extra view confidence, at the rate given by
\eqref{eq:dceff-exact}. The effective confidence exceeds the stated
confidence. For weakly correlated,
high-variance pairs or negligible $\alpha$, $c_{\eff}$ stays close to
$c$ and the estimator switch is harmless.

\subsection{From $\Delta c_{\eff}$ to $\Delta\mu$ and $\Delta w$}
\label{sec:dmu-dw}

The shifts $\Delta\mu$ and $\Delta w$ follow analytically from
$\Delta c_{\eff}$. The key step is an exact closed form for the
posterior under a single relative view, which we derive directly from
the posterior equation.

With a single view, $P = p^\top$ and $\Omega = \omega$ is a scalar; the
BL posterior equation reads
\begin{align*}
  M\,\muBL \;=\; (\tau\Sigma)^{-1}\pi \;+\; \tfrac{1}{\omega}\,p\,Q,
  \qquad M \;=\; (\tau\Sigma)^{-1} \;+\; \tfrac{1}{\omega}\,p\,p^\top.
\end{align*}
Left-multiplication by $\tau\Sigma$ and solving for $\muBL$:
\begin{align*}
  \tau\Sigma\,M\,\muBL \;&=\; \tau\Sigma\,\bigl[(\tau\Sigma)^{-1}\pi + \tfrac{1}{\omega}\,p\,Q\bigr]\\
  \Bigl[I \;+\; \tfrac{\tau}{\omega}\,\Sigma p\,p^\top\Bigr]\,\muBL \;&=\; \pi \;+\; \tfrac{\tau Q}{\omega}\,\Sigma p\\
  \muBL \;+\; \tfrac{\tau\,(p^\top\muBL)}{\omega}\,\Sigma p \;&=\; \pi \;+\; \tfrac{\tau Q}{\omega}\,\Sigma p\\
  \muBL \;&=\; \pi \;+\; \tfrac{\tau\,(Q - p^\top\muBL)}{\omega}\,\Sigma p.
  \tag{$\star$}
\end{align*}

Applying $p^\top$ to $(\star)$ produces a scalar equation in
$p^\top\muBL$ with $v_k := p^\top\Sigma p$:
\begin{align*}
  p^\top\muBL \;&=\; p^\top\pi \;+\; \tfrac{\tau v_k\,(Q - p^\top\muBL)}{\omega}\\
  \omega\,p^\top\muBL \;&=\; \omega\,p^\top\pi \;+\; \tau v_k\,Q \;-\; \tau v_k\,p^\top\muBL\\
  (\omega + \tau v_k)\,p^\top\muBL \;&=\; \omega\,p^\top\pi \;+\; \tau v_k\,Q\\
  p^\top\muBL \;&=\; \tfrac{\omega}{\omega + \tau v_k}\,p^\top\pi \;+\; \tfrac{\tau v_k}{\omega + \tau v_k}\,Q\\
  \;&=\; (1 - c_{\eff})\,p^\top\pi \;+\; c_{\eff}\,Q,
  \qquad c_{\eff} := \tfrac{\tau v_k}{\tau v_k + \omega}.
\end{align*}

Hence $Q - p^\top\muBL = (1-c_{\eff})(Q - p^\top\pi)$. Inserting this
into $(\star)$ and simplifying the prefactor:
\begin{align*}
  \muBL \;&=\; \pi \;+\; \tfrac{\tau\,(1 - c_{\eff})}{\omega}\,(Q - p^\top\pi)\,\Sigma p\\
  \;&=\; \pi \;+\; \tfrac{\tau}{\omega + \tau v_k}\,(Q - p^\top\pi)\,\Sigma p\\
  \;&=\; \pi \;+\; c_{\eff}\,(Q - p^\top\pi)\,\tfrac{\Sigma p}{v_k}.
\end{align*}

The exact identity is established:
\begin{equation}
\label{eq:bl-single-exact}
  \muBL \;=\; \pi \;+\; c_{\eff}\,(Q - p^\top\pi)\,\frac{\Sigma p}{v_k},
  \qquad c_{\eff} \;=\; \frac{\tau v_k}{\tau v_k + \omega}.
\end{equation}
The posterior moves exactly along $\Sigma p / v_k$, normalised so that
$p^\top(\Sigma p / v_k) = 1$. Substituting \eqref{eq:bl-single-exact}
into $w_{\mathrm{BL}} = \tfrac{1}{\delta}\Sigma^{-1}\muBL$ simplifies
the weight expression as well,
\begin{equation}
\label{eq:bl-single-w}
  w_{\mathrm{BL}} \;=\; w_{\mathrm{eq}} \;+\; c_{\eff}\,(Q - p^\top\pi)\,\frac{p}{\delta\,v_k},
  \qquad w_{\mathrm{eq}} = \tfrac{1}{\delta}\Sigma^{-1}\pi.
\end{equation}
The view-driven part of $w_{\mathrm{BL}}$ lies exactly in direction
$p$, without any assumption on the volatilities.
Equation~\eqref{eq:bl-single-w} is the single-view case of the
unconstrained-weights result of \citet{HeLitterman1999} (posterior
weights are the market weights plus a position in the view
portfolio) and serves here as the carrier that propagates
$\Delta c_{\eff}$.

\paragraph{Step 1: $\Delta\mu$ from $\Delta c_{\eff}$.}
Projecting \eqref{eq:bl-single-exact} on $p$ gives
\[
  p^\top \Delta\muBL \;=\; \Delta c_{\eff}\,(Q - p^\top\pi),
\]
exact for any pair of estimators, since $p^\top(\Sigma p / v_k) = 1$
under each. The full vector shift additionally moves the direction
$\Sigma p / v_k$; holding $\Sigma$ fixed, it reduces to
\[
  \Delta\muBL \;=\; \Delta c_{\eff}\,(Q - p^\top\pi)\,\tfrac{\Sigma p}{v_k}.
\]
For $p = e_i - e_j$ with $\sigma_i \approx \sigma_j$ one has
$\Sigma p \approx (v_k/p^\top p)\,p$, hence
\begin{equation}
\label{eq:dmu-vec}
  \Delta\muBL \;\approx\; \frac{\Delta c_{\eff}\,(Q - p^\top\pi)}{p^\top p}\,p,
  \qquad
  \Delta\mu_i \;\approx\; -\Delta\mu_j \;\approx\; \tfrac{1}{2}\,\Delta c_{\eff}\,(Q - p^\top\pi).
\end{equation}
This symmetric per-component split rests entirely on
$\sigma_i \approx \sigma_j$. If the volatilities differ, the direction
$\Sigma p / v_k$ distributes the return shift asymmetrically across the
two assets and the clean antisymmetric split no longer holds; the exact
vector identity \eqref{eq:bl-single-exact} is unaffected, only its
symmetric reading depends on the assumption. The numerical example below
uses a pair with nearly identical volatilities and confirms that the
symmetric split is accurate in the homogeneous case.

\paragraph{Step 2: $\Delta w$ from $\Delta\mu$.}
Differentiating \eqref{eq:bl-single-w} with respect to $c_{\eff}$ at
fixed $\Sigma$ yields directly
\begin{equation}
\label{eq:dw-dc}
  \Delta w_{\mathrm{BL}} \;=\; \frac{\Delta c_{\eff}\,(Q - p^\top\pi)}{\delta\,v_k}\,p.
\end{equation}
The shift sits exactly along $p$ with per-component
$\Delta w_i = -\Delta w_j = \Delta c_{\eff}\,(Q-p^\top\pi)/(\delta v_k)$.
\eqref{eq:dw-dc} is the exact partial derivative at fixed $\Sigma$,
i.e.\ the pure view-absorption channel. The $2\times 2$ example below,
which uses $\hat\Sigma_s^{-1}$ for both scenarios, measures it up to
the small direction shift
$\Sigma_{\LW}p/v_k^{\LW} - \hat\Sigma_s p/v_k^{\Std}$ that the shared
inverse does not remove. The separate effect of a changed $\Sigma^{-1}$
is the precision channel from Section~\ref{sec:precision}.

Every percentage point of $\Delta c_{\eff}$ moves the spread
$w_i - w_j$ linearly, with leverage $1/(\delta v_k)$. Small $v_k$
amplifies the effect, and small $v_k$ is exactly the regime in which
$\Delta c_{\eff}$ is large. Affected pairs therefore combine the largest
confidence drift with the strongest translation factor into the weights.

\subsection{A Worked $2\times 2$ Example: IWF and IWD}
\label{sec:example}

We isolate the effect on the expected-return vector by reducing the
universe to two assets, IWF (US large-cap growth) and IWD (US large-cap
value), and running the BL formula in closed form. No other assets
interact, so every number below is fully traceable. All quantities are
annualised and refer to the estimation window January 2008 to December
2018. With $\sigma_{\mathrm{IWF}} = 15.5\,\%$ and
$\sigma_{\mathrm{IWD}} = 15.6\,\%$ ($k = \sigma_i/\sigma_j = 0.99$) the
pair lies in the homogeneous special case, so the symmetric
approximations of Section~\ref{sec:dmu-dw} apply.

To keep the example self-contained, we shrink the extracted $2\times 2$
sample block with its own target
$\mu_s^{(2)} = \tfrac{1}{2}\mathrm{tr}(\hat\Sigma_s^{(2)})$, while the
intensity $\alpha = 0.041$ is the analytic value from the full
eight-asset universe. (If we shrink the full $8\times 8$ matrix instead
and extract the pair afterwards, the design underlying
Table~\ref{tab:pairs}, the larger universe-wide target
$\mu_s = 0.0282 > \mu_s^{(2)} = 0.0242$ widens the spread variance
further and the shift grows to $+10.7$ percentage points; the
self-contained design is thus the conservative one.)

The inputs are
\begin{align*}
    \delta = 2.44, \qquad \tau = 0.05, \qquad c = 0.50, \qquad
    w_{\mathrm{eq}} = (0.50,\,0.50)^\top, \qquad Q = 0.03.
\end{align*}
The two annualised covariance matrices read
\begin{align*}
    \hat{\Sigma}_s &= \begin{pmatrix} 0.02408 & 0.02227 \\ 0.02227 & 0.02435 \end{pmatrix},
    \qquad \rho_s = 0.920,\\
    \hat{\Sigma}_{\LW} &= \begin{pmatrix} 0.02409 & 0.02137 \\ 0.02137 & 0.02435 \end{pmatrix},
    \qquad \rho_{\LW} = 0.882, \qquad \alpha = 0.041.
\end{align*}
We calibrate $\omega$ once on the sample estimator and freeze it,
$\omega = \tfrac{1-c}{c}\,\tau\,v_k^\Std = 1.95 \cdot 10^{-4}$.

The effective view absorption follows directly,
\begin{align*}
    v_k^\Std &= 0.00389, & v_k^\LW &= 0.00570,\\
    c_{\eff}(\hat{\Sigma}_s) &= 50.00\,\%, & c_{\eff}(\hat{\Sigma}_{\LW}) &= 59.42\,\%,\\
    & & \Delta c_{\eff} &= +9.42\,\text{pp}.
\end{align*}

The expected-return vector and the corresponding portfolio weights shift
accordingly. The LW weights deliberately use $\hat{\Sigma}_s^{-1}$ as
well, so only the view-absorption channel shows up; the separate
precision channel through $\hat{\Sigma}^{-1}$ is treated in
Section~\ref{sec:precision}.
\begin{align*}
    \mu_{\mathrm{BL}}^\Std &= (\,6.37\,\%,\ 4.89\,\%\,), & \mu_{\mathrm{BL}}^\LW &= (\,6.52\,\%,\ 4.75\,\%\,),\\
    w_{\mathrm{BL}}^\Std &= (\,+209.4\,\%,\ -109.4\,\%\,), & w_{\mathrm{BL}}^\LW &= (\,+239.6\,\%,\ -139.2\,\%\,).
\end{align*}
The difference vectors read
\[
    \boxed{\;
    \Delta c_{\eff} = +9.42\,\text{pp}, \qquad
    \Delta\mu = (+0.15\,\text{pp},\ -0.13\,\text{pp}), \qquad
    \Delta w = (+30.2\,\text{pp},\ -29.8\,\text{pp}).
    \;}
\]

The closed-form formulas from the previous subsection reproduce these
numbers. With $p^\top p = 2$, $Q - p^\top\pi = 0.0303$, $\delta = 2.44$
and $v_k = 0.00389$ from the sample estimator we check
\begin{align*}
    \Delta\mu_{\mathrm{IWF}} \;&\approx\; \tfrac{1}{2}\,\Delta c_{\eff}\,(Q-p^\top\pi)
    \;=\; \tfrac{1}{2}\cdot 0.0942 \cdot 0.0303
    \;=\; +0.143\,\text{pp},\\
    \Delta w_{\mathrm{IWF}} \;&\approx\; \frac{\Delta c_{\eff}\,(Q-p^\top\pi)}{\delta\,v_k}
    \;=\; \frac{0.0942 \cdot 0.0303}{2.44 \cdot 0.00389}
    \;\approx\; +0.300
    \;=\; +30.0\,\text{pp}.
\end{align*}
The IWD counterparts have the opposite sign. The $\Delta\mu$ prediction
matches the numerical value to within $0.01$\,pp, the $\Delta w$
prediction to within $0.2$\,pp; both gaps stem from the slight
asymmetry $\sigma_{\mathrm{IWF}} \neq \sigma_{\mathrm{IWD}}$ that the
approximation $\Sigma p / v_k \approx p / (p^\top p)$ ignores.

The chain runs from the correlation drop to the weights: the
correlation between IWF and IWD
falls from $0.920$ to $0.882$ under shrinkage, $v_k$ grows from
$0.00389$ to $0.00570$, and $c_{\eff}$ climbs by $9.42$ percentage
points. The posterior pushes IWF up by $0.15$ pp and IWD down by $0.13$
pp. Through $w = \tfrac{1}{\delta}\Sigma^{-1}\mu$ this return shift
translates into roughly $\pm 30$ percentage points on the long-short
spread. The analyst's input form still shows $c = 50\,\%$; the model
operates at $c_{\eff} = 59.42\,\%$. The absolute position sizes
illustrate a known property of unconstrained Black--Litterman with
highly correlated pairs: the translation factor
$1/(\delta v_k) \approx 105$ turns any view into triple-digit tilts,
and the leakage is amplified by the same factor.

\subsection{Implementation in Pseudocode}
\label{sec:pseudocode}

The whole computation reduces to a few lines. Inputs are the sample
covariance, the analytic Ledoit--Wolf intensity, the market weights,
$\delta$, $\tau$, and a single relative view $(p, Q, c)$. The function
\texttt{mu\_bl} solves the Black--Litterman posterior linear system;
weights are computed with $\hat\Sigma_s^{-1}$ in both scenarios to
isolate the view-absorption channel.

\begin{lstlisting}[caption={Single-view leakage computation.},label={lst:pseudo}]
# ----- Inputs -----
Sigma_s, alpha, w_eq, delta, tau    # market and estimator parameters
p, Q, c                             # single relative view: p = e_i - e_j

# ----- Ledoit-Wolf estimator -----
n        = dim(Sigma_s)
mu_s     = trace(Sigma_s) / n
Sigma_lw = (1 - alpha) * Sigma_s + alpha * mu_s * I_n

# ----- View variance under both estimators -----
v_std = p.T @ Sigma_s  @ p
v_lw  = p.T @ Sigma_lw @ p

# ----- Idzorek omega, frozen on Sample-Sigma -----
omega = ((1 - c) / c) * tau * v_std

# ----- Effective view absorption -----
c_eff_std = tau * v_std / (tau * v_std + omega)   # equals c by construction
c_eff_lw  = tau * v_lw  / (tau * v_lw  + omega)

# ----- BL posterior (single relative view) -----
pi = delta * Sigma_s @ w_eq

def mu_bl(Sigma):
    M = inv(tau * Sigma) + outer(p, p) / omega
    b = inv(tau * Sigma) @ pi + p * Q / omega
    return solve(M, b)

mu_BL_std = mu_bl(Sigma_s)
mu_BL_lw  = mu_bl(Sigma_lw)

# ----- Weights (return-channel isolation: same Sigma^-1) -----
w_BL_std = solve(Sigma_s, mu_BL_std) / delta
w_BL_lw  = solve(Sigma_s, mu_BL_lw)  / delta

# ----- Output -----
delta_c_eff = c_eff_lw - c_eff_std      # scalar
delta_mu    = mu_BL_lw - mu_BL_std      # (n,) vector
delta_w     = w_BL_lw  - w_BL_std       # (n,) vector
\end{lstlisting}

\noindent
On the IWF/IWD inputs from Section~\ref{sec:example} the routine returns
\texttt{delta\_c\_eff = +0.0942}, \texttt{delta\_mu = (+0.0015, -0.0013)}
and \texttt{delta\_w = (+0.302, -0.298)}, matching the boxed result.

\subsection{Recommendation}
\label{sec:rec}

The clean fix is to recompute $\omega_k$ alongside $\Sigma$,
\[
  \omega_k^{\LW} \;=\; \frac{1-c_k}{c_k}\,\tau\,v_k^{\LW}.
\]
This restores $c_{\eff,k} = c_k$ exactly. The same self-healing holds
for any $\Omega$ tied functionally to $\Sigma$, for instance the
He--Litterman convention $\Omega = \mathrm{diag}(P\,\tau\Sigma\,P^\top)$:
the leakage is a property of frozen scalars, independent of the
particular calibration formula. Covariance estimator, confidence and
view-uncertainty form one unit. Replacing one of the three pieces
without re-deriving the other two is the source of the leakage, not the
choice of estimator itself. An analyst who freezes $\Omega$ has raised
the confidence of every affected view without a decision to do so, in
conflict with the risk-governance process that produced the original
confidences.

One caveat delimits the recommendation. It presumes that the analyst's
primitive is the stated confidence $c_k$, as it is under Idzorek's
procedure and under any sign-off process phrased in confidences. If the
primitive is instead $\omega_k$ itself, an absolute variance attached
to the view noise irrespective of the covariance model, then freezing
$\omega_k$ is Bayesianly coherent \citep[cf.][]{KolmRitter2017}, and
the induced change in $c_{\eff,k}$ is the correct posterior response to
a better estimate of $v_k$; recomputing $\Omega$ would then be the
error. The two readings prescribe opposite fixes, which is why the
choice of primitive belongs in the model documentation. Throughout this
paper we take $c_k$ as the primitive.

% ============================================================
\section{Conditioning and Precision Channel}
\label{sec:precision}
% ============================================================

The Idzorek case showed that the view absorption stays stable. The
fixed-$\Omega$ case showed that the shift in the return signal becomes
material under identifiable conditions. The larger effect of the LW
switch lies in the inverse
$\hat\Sigma^{-1}$ entering
$w_{\mathrm{BL}} = \frac{1}{\delta}\Sigma^{-1}\mu_{\mathrm{BL}}$. The
weight change decomposes into two channels,
\[
    \Delta w_{\mathrm{BL}}
    \;=\; \underbrace{\frac{1}{\delta}\hat{\Sigma}_{\LW}^{-1}
          \bigl(\mu_{\mathrm{BL}}^{\LW} - \mu_{\mathrm{BL}}^{\Std}\bigr)
         }_{\text{return channel}}
    \;+\; \underbrace{\frac{1}{\delta}
          \bigl(\hat{\Sigma}_{\LW}^{-1} - \hat{\Sigma}_s^{-1}\bigr)\,
          \mu_{\mathrm{BL}}^{\Std}
         }_{\text{precision channel}}.
\]
The return channel is small because Idzorek fixes the view absorption
at $c$; the remaining terms, the prior shift and the projection shift,
are both $O(\alpha)$. Under the three-view set of
Section~\ref{sec:blm-brief} the largest change in any posterior mean
across the two estimators is $0.42\,\%$ p.a., of which the prior
channel contributes at most $0.21\,\%$
(Section~\ref{sec:idzorek-cancel}); the remainder enters through the
$\Sigma$-dependent view projection at nearly unchanged absorption
(Remark~\ref{rem:multiview}). The precision
channel is driven by the difference of inverses.

LW lifts every eigenvalue $\lambda_k$ of $\hat\Sigma_s$,
\[
    \lambda_k^{\LW} \;=\; (1-\alpha)\,\lambda_k \;+\; \alpha\,\mu_s,
    \qquad k = 1,\dots,n.
\]
The smallest eigenvalue rises by the largest factor, so $1/\lambda_n$
falls the most; on our data it falls by roughly a factor of three. As a
result $\hat\Sigma_{\LW}^{-1}$ is much more uniform than
$\hat\Sigma_s^{-1}$, and the LW inverse pulls back the extreme
long-short positions that the badly conditioned estimator produces
along low-variance directions. On our universe
the condition number drops from $341.8$ to $108.0$; under the three-view
set of Section~\ref{sec:blm-brief} (weights via
$w = \tfrac{1}{\delta}\Sigma^{-1}\muBL$), the IWF position (US
large-cap growth)
falls from $+173.5\,\%$ to $+121.0\,\%$ and the IWD short (US large-cap
value) from $-136.4\,\%$ to $-83.0\,\%$, taking the gross exposure from
$4.31$ to $3.16$.

The cost of the frozen-$\Omega$ workflow appears at the same level,
and it requires no data beyond the estimation window. With $\Omega$
frozen on the sample estimator, the inflated view confidences partially
restore the trimmed positions (IWF $+140.5\,\%$ instead of
$+121.0\,\%$; gross exposure $3.63$ instead of $3.16$) and undo roughly
$40\,\%$ of the exposure reduction that the recomputed-$\Omega$ switch
delivers. Economically, the inflated confidence acts as leverage on the
view: the position scales with $c_{\eff}$, and with it both the payoff
if the view is right and the loss if it is wrong. A performance
evaluation on subsequent data would add only a single realised path of
this trade-off; we therefore report the weight-level cost.

% ============================================================
% The out-of-sample backtest of earlier drafts is parked below in an
% \iffalse block. Rationale: the closed forms are identities, the
% weight- and exposure-level consequences of the leakage are
% estimation-window facts (now reported in Section 7), and a
% performance evaluation adds only a single realised path of the
% confidence-as-leverage trade-off. Re-enable by deleting \iffalse/\fi
% and restoring the Memmel/LedoitWolf2008 references in the
% bibliography and section [6b] of code/paper_numbers.py.
% ============================================================
\iffalse
% ============================================================
\section{Out-of-Sample Backtest}
\label{sec:backtest}
% ============================================================

All weights are estimated once on the training window (January 2008 to
December 2018) and held constant on the test window (January 2019 to
December 2024), i.e.\ the portfolios are rebalanced back to fixed
weights each month. The three relative views of
Section~\ref{sec:blm-brief} apply: IWF $>$ IWD by $+3\,\%$ p.a., EEM $>$
EFA by $+4\,\%$ p.a., IWF $>$ AGG by $+6\,\%$ p.a., each at confidence
$c = 0.50$, $\tau = 0.05$. Three Black--Litterman variants are run: the
sample estimator; Ledoit--Wolf with Idzorek-$\Omega$ recomputed per
estimator, so the view channel is clean up to the small multi-view
cross-terms of Remark~\ref{rem:multiview}; and Ledoit--Wolf with
$\Omega$ frozen on the sample estimator, the leakage workflow of
Section~\ref{sec:fixed-omega}. The view set follows Idzorek's examples
(relative pair views with moderate spreads). Because the test window is
known to us, the performance levels below are illustrative; the object
of interest is the comparison across the three $\Omega$-treatments,
which share identical views. Ex post, two views turn out
right and one wrong (realised spreads 2019--2024: IWF$-$IWD $+9.1$,
IWF$-$AGG $+20.3$, EEM$-$EFA $-3.6\,\%$ p.a.). All performance figures
are arithmetic annualisations of monthly excess returns; the maximum
drawdown is computed on the compounded excess-return path. Transaction
costs, shorting fees and margin constraints are ignored, which flatters
the levered variants; financing at the T-bill rate is implicit in the
excess-return arithmetic. Relative views leave all portfolios exactly
fully invested under the weight mapping of Section~\ref{sec:blm-brief}.

\begin{table}[H]
\centering
\renewcommand{\arraystretch}{1.15}
\caption{Out-of-sample performance (Jan.\,2019--Dec.\,2024), monthly
excess returns, weights fixed at inception (monthly rebalancing).
$\kappa(\hat\Sigma_s) = 341.8$, $\kappa(\hat\Sigma_{\LW}) = 108.0$.
The Sharpe difference between the sample and recomputed-$\Omega$ LW
columns is statistically indistinguishable (Memmel test, $p = 0.62$).}
\label{tab:backtest}
\small
\begin{tabular}{lrrrr}
\toprule
Metric & Market eq. & BLM (sample) & BLM (LW, $\Om$ rec.) & BLM (LW, $\Om$ frozen) \\
\midrule
Annual return     & $3.85\,\%$   & $18.06\,\%$  & $13.54\,\%$  & $15.09\,\%$ \\
Annual volatility & $12.16\,\%$  & $24.98\,\%$  & $19.81\,\%$  & $21.65\,\%$ \\
Sharpe ratio      & $0.32$       & $0.72$       & $0.68$       & $0.70$ \\
Max drawdown      & $-24.78\,\%$ & $-52.34\,\%$ & $-43.06\,\%$ & $-46.79\,\%$ \\
\bottomrule
\end{tabular}
\end{table}

Compare the first two BLM columns. The largest change in any posterior
mean across the two estimators is $0.42\,\%$ p.a., of which the pure
prior channel contributes at most $0.21\,\%$
(Section~\ref{sec:idzorek-cancel}); the remainder enters through the
$\Sigma$-dependent view projection at nearly unchanged absorption
(Remark~\ref{rem:multiview}). The Idzorek cancellation works as
expected, and the performance gap is generated almost entirely by the
precision channel. $\hat\Sigma_s^{-1}$ produces $+173.5\,\%$ in US
large-cap growth and $-136.4\,\%$ in US large-cap value;
$\hat\Sigma_{\LW}^{-1}$ trims these to $+121.0\,\%$ and $-83.0\,\%$.
Because growth stocks strongly outperformed in the test window, the
concentrated sample-covariance bet pays off ($18.06\,\%$, Sharpe
$0.72$). Ledoit--Wolf cuts volatility to $19.81\,\%$ and the maximum
drawdown to $-43.06\,\%$, at the cost of a portion of the upside. The
Sharpe gap is well inside sampling noise over 72 months
(\citealp{Memmel2003}; see also \citealp{LedoitWolf2008}). Fewer extreme
positions mean less risk and less return in an up-market. With Idzorek
$\Omega$ the LW switch does not change the return signal, it only
dampens the leverage of the optimisation step.

The fourth column quantifies the leakage of Section~\ref{sec:fixed-omega}
directly: $\Omega$ is calibrated on $\hat\Sigma_s$, frozen, and the
covariance is switched to LW. The silently inflated view confidences
partially restore the concentrated positions that the estimator switch
was meant to trim (IWF $+140.5\,\%$ against $+121.0\,\%$ with recomputed
$\Omega$; gross exposure $3.63$ against $3.16$), and with them roughly
$40\,\%$ of both the exposure reduction and the drawdown improvement
that the clean LW switch delivers.

LW underperforms the sample-covariance BL on annualised return and
Sharpe in this specific test window, a strong US growth market. The
mechanism is structural: LW caps the concentrated growth position, and
because US large-cap growth dominated every other asset class in the
test period, LW gives up return. The maximum drawdown, however, falls
from $-52.34\,\%$ to $-43.06\,\%$, a nine-percentage-point reduction.
The drawdown reduction is where LW adds value in this window; drawdowns
above $50\,\%$ breach common institutional loss limits. Imposing a
long-only constraint instead compresses all three variants toward each
other and toward far tamer risk figures: annualised volatilities of
$15.1\,\%$, $15.4\,\%$ and $15.7\,\%$ and maximum drawdowns of
$-28.3\,\%$, $-29.2\,\%$ and $-29.2\,\%$ for the sample,
recomputed-$\Omega$ and frozen-$\Omega$ variants (reproduced by the
accompanying code). This is consistent
with the observation of \citet[p.~1652]{JagannathanMa2003} that
no-short-sale constraints act on the same estimation-error leverage that
shrinkage targets; investors who optimise long-only to begin with gain
correspondingly little from the estimator switch.

Three configurations plausibly favour LW; we state them as
conjectures. Regulated
portfolios
with explicit drawdown or VaR constraints benefit, because the loss
limit, not the risk-adjusted return, is the binding constraint. Periods
of high market volatility or correlation regime shifts are a second
case, because badly conditioned sample estimators tend to produce
spurious extreme positions, especially after short training windows or
at low $T/n$. And third, whenever the training window contains a
structural regime anomaly (such as the unusually wide eigenvalue spreads
of 2008--2009) that enters the sample covariance and distorts the
weights held during the test window.
\fi

% ============================================================
\section{Conclusion}
\label{sec:conclusion}
% ============================================================

Within the view channel, whether the Ledoit--Wolf switch is harmless
inside Black--Litterman is
decided by the treatment of $\Omega$, and both regimes now have closed
forms. Under Idzorek's calibration
$\omega_k(\Sigma) = \tfrac{1-c_k}{c_k}\tau\,p_k^\top\Sigma p_k$ the
absorption $c_{\eff,k}$ is independent of $\Sigma$ for a single view:
any positive-definite estimator can be plugged in without distorting the
view channel, and the posterior impact is confined to prior and
projection shifts, the former exactly linear and the latter first-order
linear in $\alpha$. If
$\omega_k$ is instead computed once and frozen, the estimator switch
shifts the effective confidence by the exact, $\tau$-free expression
\eqref{eq:dceff-exact}. The sign criterion is
$v_k^\Std \lessgtr \mu_s\lVert p_k\rVert^2$ ($2\mu_s$ for relative
pairs), the shift is confined to the band
$(-c(1-c)\alpha/(1-c\alpha),\,1-c)$ of Remark~\ref{rem:range}, the
symmetric special case yields the closed-form correlation threshold
\eqref{eq:rho-star}, and the exact
single-view identity \eqref{eq:bl-single-exact} propagates the shift
into posterior means and weights with leverage $1/(\delta v_k)$. On the
eight-asset universe at $\alpha = 0.041$, five of 28 pairs cross the
five-percentage-point band, led by IWF/IWD at $+10.7$; the bond pair
AGG/BWX crosses it at a correlation of only $0.57$. In the weights, the
frozen-$\Omega$ workflow restores roughly $40\,\%$ of the
gross-exposure reduction that the estimator switch delivers.
Recomputing $\omega_k$ alongside $\Sigma$ removes the leakage
analytically. The lesson is procedural: covariance, confidence and view
uncertainty have to be kept consistent with each other.

Several limitations point to extensions. The closed-form chain covers a
single view (arbitrary view vectors in the shift formula, relative
pairs for the per-component weight split) and a diagonal $\Omega$; with
several simultaneous
views the shifts interact through the posterior precision
(Remark~\ref{rem:multiview} bounds the size of these cross-terms on our
data), and the multi-view analogue (via Sherman--Morrison updates) we
leave to future work. The empirical part
rests on a single estimation window and a small, well-diversified ETF
universe on which the analytic intensity is modest; equity universes
with larger $n/T$ push $\alpha$, and with it the leakage, substantially
higher. Finally, the mechanism is target-independent (replacing the
scaled identity by the single-factor target of \citet{LedoitWolf2003} or
the nonlinear shrinkage of \citet{LedoitWolf2017} only changes the
expression for $v_k^{\LW} - v_k^\Std$), but the quantitative reach of
the effect under these estimators remains to be mapped out.

% ============================================================
\section*{Reproducibility}
% ============================================================

All numbers quoted in this paper are produced by the script
\texttt{code/paper\_numbers.py} from frozen monthly return data
(\texttt{data/}, sourced from Yahoo Finance). Pipeline: excess returns;
sample covariance with $T-1$ normalisation; the analytic Ledoit--Wolf
intensity of \texttt{scikit-learn} applied to that matrix
(Section~\ref{sec:lw}); weights via
$w = (\delta\Sigma)^{-1}\mu_{\mathrm{BL}}$; implied risk aversion from
S\&P~500 excess returns. Code and data are available at
\url{https://github.com/lucasposern/blm-lw-fixed-omega}.

% ============================================================
\bibliographystyle{plainnat}
% Hand-maintained bibliography for the draft; switch to BibTeX with the
% target journal's style at submission.
\begin{thebibliography}{29}

\bibitem[Bessler et~al.(2017)]{BesslerOpferWolff2017}
W.~Bessler, H.~Opfer, and D.~Wolff.
\newblock Multi-asset portfolio optimization and out-of-sample
performance: An evaluation of Black--Litterman, mean-variance, and
na\"ive diversification approaches.
\newblock \emph{The European Journal of Finance}, 23(1):1--30, 2017.

\bibitem[Best and Grauer(1991)]{BestGrauer1991}
M.~J. Best and R.~R. Grauer.
\newblock On the sensitivity of mean-variance-efficient portfolios to
changes in asset means: Some analytical and computational results.
\newblock \emph{The Review of Financial Studies}, 4(2):315--342, 1991.

\bibitem[Bevan and Winkelmann(1998)]{BevanWinkelmann1998}
A.~Bevan and K.~Winkelmann.
\newblock Using the Black--Litterman global asset allocation model:
Three years of practical experience.
\newblock Fixed Income Research, Goldman Sachs \& Co., 1998.

\bibitem[Black and Litterman(1992)]{BlackLitterman1992}
F.~Black and R.~Litterman.
\newblock Global portfolio optimization.
\newblock \emph{Financial Analysts Journal}, 48(5):28--43, 1992.

\bibitem[Chopra and Ziemba(1993)]{ChopraZiemba1993}
V.~K. Chopra and W.~T. Ziemba.
\newblock The effect of errors in means, variances, and covariances on
optimal portfolio choice.
\newblock \emph{The Journal of Portfolio Management}, 19(2):6--11, 1993.

\bibitem[DeMiguel et~al.(2009)]{DeMiguelGarlappiUppal2009}
V.~DeMiguel, L.~Garlappi, and R.~Uppal.
\newblock Optimal versus naive diversification: How inefficient is the
$1/N$ portfolio strategy?
\newblock \emph{The Review of Financial Studies}, 22(5):1915--1953, 2009.

\bibitem[DeMiguel et~al.(2013)]{DeMiguelMartinUtreraNogales2013}
V.~DeMiguel, A.~Mart\'in-Utrera, and F.~J. Nogales.
\newblock Size matters: Optimal calibration of shrinkage estimators for
portfolio selection.
\newblock \emph{Journal of Banking \& Finance}, 37(8):3018--3034, 2013.

\bibitem[Frost and Savarino(1986)]{FrostSavarino1986}
P.~A. Frost and J.~E. Savarino.
\newblock An empirical Bayes approach to efficient portfolio selection.
\newblock \emph{Journal of Financial and Quantitative Analysis},
21(3):293--305, 1986.

\bibitem[He and Litterman(1999)]{HeLitterman1999}
G.~He and R.~Litterman.
\newblock The intuition behind Black--Litterman model portfolios.
\newblock Investment Management Research, Goldman Sachs \& Co., 1999.

\bibitem[Herold(2003)]{Herold2003}
U.~Herold.
\newblock Portfolio construction with qualitative forecasts.
\newblock \emph{The Journal of Portfolio Management}, 30(1):61--72, 2003.

\bibitem[Idzorek(2007)]{Idzorek2007}
T.~Idzorek.
\newblock A step-by-step guide to the Black--Litterman model:
Incorporating user-specified confidence levels.
\newblock In S.~Satchell, editor, \emph{Forecasting Expected Returns in
the Financial Markets}, pages 17--38. Academic Press, 2007.

\bibitem[Jagannathan and Ma(2003)]{JagannathanMa2003}
R.~Jagannathan and T.~Ma.
\newblock Risk reduction in large portfolios: Why imposing the wrong
constraints helps.
\newblock \emph{The Journal of Finance}, 58(4):1651--1683, 2003.

\bibitem[Jorion(1986)]{Jorion1986}
P.~Jorion.
\newblock Bayes--Stein estimation for portfolio analysis.
\newblock \emph{Journal of Financial and Quantitative Analysis},
21(3):279--292, 1986.

\bibitem[Kan and Zhou(2007)]{KanZhou2007}
R.~Kan and G.~Zhou.
\newblock Optimal portfolio choice with parameter uncertainty.
\newblock \emph{Journal of Financial and Quantitative Analysis},
42(3):621--656, 2007.

\bibitem[Kolm and Ritter(2017)]{KolmRitter2017}
P.~N. Kolm and G.~Ritter.
\newblock On the Bayesian interpretation of Black--Litterman.
\newblock \emph{European Journal of Operational Research},
258(2):564--572, 2017.

\bibitem[Ledoit and Wolf(2003)]{LedoitWolf2003}
O.~Ledoit and M.~Wolf.
\newblock Improved estimation of the covariance matrix of stock returns
with an application to portfolio selection.
\newblock \emph{Journal of Empirical Finance}, 10(5):603--621, 2003.

\bibitem[Ledoit and Wolf(2004{\natexlab{a}})]{LedoitWolf2004}
O.~Ledoit and M.~Wolf.
\newblock A well-conditioned estimator for large-dimensional covariance
matrices.
\newblock \emph{Journal of Multivariate Analysis}, 88(2):365--411,
2004{\natexlab{a}}.

\bibitem[Ledoit and Wolf(2004{\natexlab{b}})]{LedoitWolfHoney2004}
O.~Ledoit and M.~Wolf.
\newblock Honey, I shrunk the sample covariance matrix.
\newblock \emph{The Journal of Portfolio Management}, 30(4):110--119,
2004{\natexlab{b}}.

% Parked with the backtest section (re-enable together):
% \bibitem[Ledoit and Wolf(2008)]{LedoitWolf2008}
% O.~Ledoit and M.~Wolf.
% \newblock Robust performance hypothesis testing with the Sharpe ratio.
% \newblock \emph{Journal of Empirical Finance}, 15(5):850--859, 2008.

\bibitem[Ledoit and Wolf(2017)]{LedoitWolf2017}
O.~Ledoit and M.~Wolf.
\newblock Nonlinear shrinkage of the covariance matrix for portfolio
selection: Markowitz meets Goldilocks.
\newblock \emph{The Review of Financial Studies}, 30(12):4349--4388, 2017.

\bibitem[Ledoit and Wolf(2020)]{LedoitWolf2020}
O.~Ledoit and M.~Wolf.
\newblock Analytical nonlinear shrinkage of large-dimensional
covariance matrices.
\newblock \emph{The Annals of Statistics}, 48(5):3043--3065, 2020.

\bibitem[Markowitz(1952)]{Markowitz1952}
H.~Markowitz.
\newblock Portfolio selection.
\newblock \emph{The Journal of Finance}, 7(1):77--91, 1952.

% Parked with the backtest section (re-enable together):
% \bibitem[Memmel(2003)]{Memmel2003}
% C.~Memmel.
% \newblock Performance hypothesis testing with the Sharpe ratio.
% \newblock \emph{Finance Letters}, 1(1):21--23, 2003.

\bibitem[Meucci(2010)]{Meucci2010}
A.~Meucci.
\newblock The Black--Litterman approach: Original model and extensions.
\newblock In R.~Cont, editor, \emph{Encyclopedia of Quantitative
Finance}. Wiley, Chichester, 2010.

\bibitem[Michaud(1989)]{Michaud1989}
R.~O. Michaud.
\newblock The Markowitz optimization enigma: Is `optimized' optimal?
\newblock \emph{Financial Analysts Journal}, 45(1):31--42, 1989.

\bibitem[Michaud et~al.(2013)]{MichaudEschMichaud2013}
R.~O. Michaud, D.~N. Esch, and R.~Michaud.
\newblock Deconstructing Black--Litterman: How to get the portfolio you
already knew you wanted.
\newblock \emph{Journal of Investment Management}, 11(1):6--20, 2013.

\bibitem[Satchell and Scowcroft(2000)]{SatchellScowcroft2000}
S.~Satchell and A.~Scowcroft.
\newblock A demystification of the Black--Litterman model: Managing
quantitative and traditional portfolio construction.
\newblock \emph{Journal of Asset Management}, 1(2):138--150, 2000.

\bibitem[Sharpe(1964)]{Sharpe1964}
W.~F. Sharpe.
\newblock Capital asset prices: A theory of market equilibrium under
conditions of risk.
\newblock \emph{The Journal of Finance}, 19(3):425--442, 1964.

\bibitem[Walters(2014)]{Walters2014}
J.~Walters.
\newblock The Black--Litterman model in detail.
\newblock Working paper, SSRN 1314585, 2014.

\end{thebibliography}

\end{document}
